\section{Discussão --- o que mudou, por que mudou e quanto se pode confiar}\label{cap:discussao}

Os capítulos anteriores descrevem \emph{o que} mudou na literatura entre os dois
períodos. Este capítulo pergunta \emph{por que} mudou e confronta os achados com
o referencial teórico do Capítulo~\ref{cap:refteo}. Ele se organiza em torno de
quatro perguntas que a comparação bruta pré/pós deixa em aberto: se o
deslocamento do risco acompanha a \emph{data} ou o \emph{objeto} tecnológico dos
estudos; se o sinal sobre o emprego depende do que cada desenho mede; quanto das
mudanças é recomposição do corpus e quanto é mudança dentro de cada tipo de
estudo; e quão sensíveis são as conclusões às escolhas metodológicas da revisão.
As análises são, como no restante do trabalho, exploratórias.

\subsection{Data ou objeto? O risco acompanha a tecnologia estudada}

A comparação pré/pós agrupa os estudos pela data de publicação, mas a hipótese
da revisão é, no fundo, uma hipótese sobre \emph{tecnologia}: a de que a IA
generativa atinge tarefas distintas das que a automação atingia. As duas coisas
não coincidem, pois apenas 9,2\% dos estudos publicados de 2023 em diante têm a
IA generativa como objeto (Capítulo~\ref{cap:descritivas}). A
Tabela~\ref{tab:polarizacao_tecnologia} separa as duas dimensões, e o resultado é
o mais nítido de toda a revisão. Entre os estudos cujo objeto é a automação
clássica ou a robótica, o risco para a alta qualificação é praticamente ausente
(2,1\% e 3,4\%) e a baixa qualificação é o diagnóstico de dois terços dos
trabalhos. Entre os estudos cujo objeto é a IA generativa, a ordem se inverte: o
risco para a alta qualificação é a categoria \emph{modal} (48,5\%), e a baixa
qualificação cai a 27,3\%. Os estudos que tratam a IA de modo geral ocupam
posição intermediária (17,0\%).

\begin{table}[htbp]
\centering
\caption{Polarização segundo a tecnologia em foco no estudo.}
\label{tab:polarizacao_tecnologia}
\input{tables/polarizacao_tecnologia.tex}
\end{table}

A leitura teórica é direta, e é favorável ao arcabouço de tarefas. O modelo de
\textcite{autorLevyMurnane2003} nunca afirmou que a tecnologia ameaça a baixa
qualificação; afirmou que ela substitui as tarefas que consegue executar, e que
essas tarefas eram, para o computador e para o robô, as rotineiras. O que os
índices de exposição da geração seguinte mostraram é que o conjunto de tarefas
executáveis pela IA é outro \parencite{webb2020, felten2021, eloundou2024}. O
corpus reproduz exatamente essa estrutura: a literatura não mudou de ideia sobre
a automação --- nos estudos sobre automação e robôs, a baixa qualificação é o
diagnóstico de 68,8\% dos trabalhos no pré e de 67,5\% no pós, e a alta
qualificação, de 1,8\% e 3,1\% ---, mas passou a estudar também uma
tecnologia cujo perfil de tarefas é diferente, e para ela o diagnóstico é
diferente. A ruptura está no \emph{objeto}; a continuidade está na
\emph{teoria}. Nesse sentido, a hipótese de trabalho da revisão é confirmada onde
ela de fato se aplica, e a expressão ``ruptura parcial'' ganha conteúdo preciso:
a ruptura é parcial porque a IA generativa ainda é o objeto de uma parcela
pequena da literatura indexada, não porque o diagnóstico sobre ela seja
hesitante.

\begin{table}[htbp]
\centering
\caption{Determinantes de o estudo situar o risco na alta qualificação (logit).}
\label{tab:logit_h1}
{\small\input{tables/logit_h1.tex}}
\end{table}

O objeto, contudo, não esgota a explicação. A Tabela~\ref{tab:logit_h1} mostra
que, controlado o objeto generativo, publicar de 2023 em diante ainda eleva as
chances de situar o risco na alta qualificação (razão de chances de 4,06 na
coluna 2 e de 3,18 com todos os controles), e a
Tabela~\ref{tab:h1_sensibilidade} mostra que o efeito do período sobrevive à
\emph{exclusão} de todos os estudos de IA generativa (2,7\% para 10,2\%). O
canal é visível nos estudos de IA ``geral'': neles, a fração que aponta a alta
qualificação passa de 1 em 32 no pré para 28 em 139 no pós. Há, portanto, um
efeito de \emph{difusão do enquadramento}: a chegada da IA generativa alterou a
forma como a literatura fala de IA em sentido amplo, mesmo quando o estudo não
trata de modelos generativos. A Figura~\ref{fig:h1_tendencia} sugere ainda que
esse movimento é gradual e anterior ao pivô: a fração cresce de forma monotônica
desde 2021 (razão de chances de 1,46 por ano, $p < 0{,}001$; 1,38 sem os estudos
de IA generativa), em linha com o fato de que os índices de exposição à IA
preditiva já apontavam para a alta qualificação antes do ChatGPT
\parencite{webb2020, felten2021}. O lançamento do ChatGPT acelera uma tendência;
não a inaugura.

\begin{figure}[htbp]
\centering
\includegraphics[width=0.85\linewidth]{figures/h1_tendencia_anual.pdf}
\caption{Fração dos estudos que situam o risco na alta qualificação, por ano de
publicação (faixa: IC Wilson 95\%; $n$ = estudos que classificaram a polarização).}
\label{fig:h1_tendencia}
\end{figure}

\begin{table}[htbp]
\centering
\caption{Sensibilidade do teste da hipótese central a definições e subconjuntos.}
\label{tab:h1_sensibilidade}
{\small\input{tables/h1_sensibilidade.tex}}
\end{table}

A Tabela~\ref{tab:h1_sensibilidade} documenta também a fragilidade que esta
revisão corrigiu. Na primeira versão da análise, a posição pré/pós era atribuída
pelo modelo de linguagem extrator, que, sem instrução explícita, classificou
como ``pré'' 115 estudos publicados entre 2023 e 2026 cujos \emph{dados} eram
anteriores a 2022. Com aquela variável o efeito era fraco (7,5\% para 12,9\%,
$p = 0{,}0609$); com o ano de publicação, que é o critério declarado na
metodologia, ele é forte. O resultado é insensível à escolha do pivô (2023 ou
2024), à exclusão dos anos de transição, à restrição a estudos revisados por
pares ou empíricos e à exclusão dos estudos sobre a China. Perde significância
apenas nos dois menores subconjuntos --- estudos de \textit{score} $\geq 4$ e
estudos lidos em texto completo ---, nos quais a direção se mantém, mas o grupo
pré tem 34 e 27 estudos classificados, respectivamente.

\subsection{O sinal depende do que se mede}

A Introdução conjecturou que os diagnósticos da literatura ``divergem menos por
erro do que por medirem objetos distintos''. A
Tabela~\ref{tab:sinal_por_tipo} submete a conjectura aos dados, e ela se
sustenta: a associação entre o tipo de evidência e o sinal sobre o emprego ($V$
de Cramér de 0,30) é a mais forte do corpus, superior à de qualquer dimensão com
o período. Metade dos estudos de exposição ocupacional conclui por efeito
negativo; mais de 40\% dos estudos de firma concluem por efeito positivo; e cerca
de 80\% dos trabalhos teóricos e das revisões declaram o efeito ambíguo.

\begin{table}[htbp]
\centering
\caption{Sinal sobre o emprego segundo o tipo de evidência.}
\label{tab:sinal_por_tipo}
\input{tables/sinal_por_tipo.tex}
\end{table}

Cada um desses padrões tem explicação no próprio arcabouço de
\textcite{acemogluRestrepo2019}. Um índice de exposição mede a viabilidade
técnica de executar tarefas por máquina; ele captura, por construção, o
efeito-deslocamento e nada diz sobre reinstalação, produtividade ou demanda. A
crítica de \textcite{arntzGregoryZierahn2017} a \textcite{frey2017} é, no fundo,
esta. Um estudo de firma, no extremo oposto, compara firmas adotantes com não
adotantes: as adotantes tornam-se mais produtivas, ganham mercado e contratam,
de modo que o desenho captura o efeito-produtividade em equilíbrio parcial. Mas o
emprego que a adotante ganha pode ser o que a concorrente perde, e o saldo
positivo da firma não se transpõe ao setor, ponto que \textcite{bessen2018}
formaliza pela elasticidade da demanda. Os modelos teóricos, por fim, são
``ambíguos'' por definição: o seu resultado é justamente o de que o saldo depende
de parâmetros. A heterogeneidade de vereditos que motivou esta revisão é, assim,
em larga medida um artefato previsível da divisão de trabalho entre os desenhos,
e não um desacordo substantivo. A evidência de nível agregado disponível é
coerente com essa leitura: efeitos nulos ou pequenos sobre o emprego total
\parencite{acemogluVacancies2022, kauhanenRouvinen2025}, compatíveis com a
estimativa de \textcite{acemoglu2025simple} de que os ganhos macroeconômicos da
IA na próxima década serão modestos, convivem com efeitos negativos rápidos e
mensuráveis onde a substituição é direta e o mercado é flexível, como no trabalho
\textit{freelance} de redação e de criação de imagens \parencite{huiReshefZhou2024}.

Decorre daí uma implicação metodológica para quem lê esta literatura: a
proporção de estudos ``negativos'' ou ``positivos'' num corpus não é um
termômetro do efeito da IA, e sim, em boa parte, um retrato da composição de
desenhos desse corpus.

\subsection{Recomposição do corpus ou mudança de diagnóstico?}

Se o sinal depende do desenho e a composição de desenhos mudou
(Capítulo~\ref{cap:comparacao}), é preciso perguntar quanto das diferenças
pré/pós é mera recomposição. A Tabela~\ref{tab:composicao} compara, para cada
desfecho, a razão de chances bruta do período com a razão ajustada por tipo de
evidência, objeto generativo, foco na China e qualidade.

\begin{table}[htbp]
\centering
\caption{Diferenças pré/pós brutas e ajustadas pela composição do corpus.}
\label{tab:composicao}
{\small\input{tables/composicao_ajuste.tex}}
\end{table}

Três padrões emergem. O crescimento do sinal \emph{positivo} é, em cerca de
metade, composição: a razão de chances cai de 5,12 para 2,49 quando se controla o
tipo de estudo, pois a literatura recente tem quatro vezes mais estudos de firma,
que são os que encontram efeitos positivos. O recuo do \emph{deslocamento} é
quase todo composição (de 0,27 para 0,52, já sem significância a 5\%). Em
contraste, a queda do sinal \emph{negativo}, a alta da complementaridade, a queda
da demanda agregada e o encurtamento do horizonte mantêm-se após o ajuste: são
mudanças que ocorrem \emph{dentro} de cada tipo de estudo. E o deslocamento do
risco para a alta qualificação, embora atenuado, permanece.

\begin{table}[htbp]
\centering
\caption{Perfis de mecanismos: participação por período e sinal associado.}
\label{tab:perfis_mecanismos}
\input{tables/perfis_mecanismos.tex}
\end{table}

A Tabela~\ref{tab:perfis_mecanismos} reorganiza os quatro mecanismos em perfis e
oferece a ponte mais direta entre a teoria e os resultados. O arcabouço de
Acemoglu e Restrepo afirma que o efeito da tecnologia sobre o emprego é o saldo
entre uma força de deslocamento e forças compensatórias. O corpus raciocina
exatamente assim: dos estudos que invocam \emph{apenas} o deslocamento, 67,2\%
concluem por efeito negativo; dos que invocam \emph{apenas} forças
compensatórias, 64,8\% concluem por efeito positivo; e dos que invocam as duas
coisas --- três quartos do corpus ---, 73,0\% declaram o efeito ambíguo. A
ambiguidade predominante da literatura não é indecisão: é a conclusão que se
segue de reconhecer as duas forças sem dispor de magnitudes comparáveis para
pesá-las, lacuna documentada na Tabela~\ref{tab:magnitude_cobertura}. O que muda
entre os períodos é o crescimento do perfil ``só compensação'', de 3,9\% para
14,1\% dos estudos.

Dois traços dessa mudança merecem leitura crítica à luz do referencial. O
primeiro é a assimetria entre complementaridade, que sobe quinze pontos, e
reinstalação, que não se move. No arcabouço teórico as duas não são equivalentes:
é a criação de \emph{novas tarefas}, e não o ganho de produtividade nas
existentes, que sustenta a demanda por trabalho no longo prazo
\parencite{acemogluRestrepo2019}. Uma literatura que documenta cada vez mais que
a IA ajuda o trabalhador a fazer melhor o que já fazia, sem documentar na mesma
proporção o surgimento de tarefas novas, descreve um cenário mais próximo da
advertência de \textcite{acemogluRestrepo2020wrongai} e da ``armadilha de
Turing'' de \textcite{brynjolfsson2022turing} do que do otimismo que a palavra
``complementaridade'' sugere: ganhos de produtividade em tarefas existentes
podem, com demanda inelástica, reduzir o emprego nessas mesmas tarefas.

O segundo traço é o recuo da demanda agregada, que cai pela metade. A literatura
torna-se mais empírica e mais microeconômica, o que é um ganho de credibilidade,
mas abandona justamente o canal do qual depende a resposta à pergunta agregada.
Forma-se um ponto cego: sabe-se cada vez mais sobre o que acontece na firma que
adota e cada vez menos sobre o que acontece na economia em que ela está.

Há, por fim, um resultado que o esquema de codificação desta revisão não
consegue capturar e que a Introdução já insinuava. A evidência experimental
mostra ganhos concentrados nos trabalhadores \emph{menos} experientes dentro de
uma mesma ocupação \parencite{noyZhang2023, brynjolfssonLiRaymond2025}. Isso não
é risco para a ``baixa'' nem para a ``alta'' qualificação entre ocupações: é
\emph{compressão} do retorno à experiência dentro delas. \textcite{autor2024rebuild}
constrói sobre essa evidência a conjectura de que a IA pode baratear a
\emph{expertise} e reabrir tarefas qualificadas a trabalhadores de qualificação
intermediária, revertendo parte da polarização; \textcite{agrawalGansGoldfarb2019}
já indicavam que automatizar a predição pode tanto substituir quanto valorizar o
julgamento humano. A dicotomia alta/baixa qualificação, herdada do debate
SBTC--RBTC, é possivelmente estreita demais para a IA generativa, e a categoria
``ambos'', estável em cerca de um terço do corpus, pode estar abrigando parte
desse fenômeno. Uma futura rodada de extração deveria codificar a dimensão
intraocupacional separadamente.

\subsection{Quanto se pode confiar: validade e sensibilidade}

\textbf{Comparações múltiplas.} O Capítulo~\ref{cap:comparacao} realiza nove
testes pré/pós. A Tabela~\ref{tab:multiplos_testes} aplica a correção de Holm: à
exceção da reinstalação, que já não era significativa, todas as associações se
mantêm.

\begin{table}[htbp]
\centering
\caption{Testes pré/pós com correção para comparações múltiplas.}
\label{tab:multiplos_testes}
\input{tables/multiplos_testes.tex}
\end{table}

\textbf{Extração a partir do resumo.} Como 86\% dos estudos foram codificados a
partir do resumo, a Tabela~\ref{tab:validade_fonte} compara-os com os lidos em
texto completo. Para sinal, tipo de evidência e horizonte não há diferença
detectável. Para a polarização há: o resumo deixa a dimensão sem classificar em
39,9\% dos casos, contra 11,4\% no texto completo, e, quando classifica,
simplifica --- a categoria ``ambos'' corresponde a 45,5\% dos estudos lidos na
íntegra e a 28,4\% dos lidos pelo resumo. A variável central da revisão é,
portanto, a mais afetada pela limitação de acesso ao texto completo, e as
proporções de ``baixa'' e de ``alta'' qualificação estão provavelmente
superestimadas em detrimento de ``ambos''. Não há, porém, indício de que o viés
difira entre os períodos.

\begin{table}[htbp]
\centering
\caption{Validade da extração: estudos lidos pelo resumo vs.\ em texto completo.}
\label{tab:validade_fonte}
\input{tables/validade_fonte.tex}
\end{table}

\textbf{Método comum.} Todas as variáveis substantivas foram codificadas pelo
mesmo modelo, a partir do mesmo texto e na mesma chamada. Associações entre elas
--- como a dos perfis de mecanismos com o sinal --- podem estar infladas pela
tendência do extrator a produzir codificações internamente coerentes. As
associações com o ano de publicação, que é determinístico, não sofrem desse
problema, e é nelas que se apoia o teste da hipótese central. A verificação
humana amostral descrita na metodologia é o instrumento para dimensionar esse
risco e permanece como a principal pendência do trabalho.

\textbf{Cobertura da busca.} A rede de citações do Capítulo~\ref{cap:bibliometria}
permite aferir a sensibilidade da busca contra um padrão externo: as referências
mais citadas pelo próprio corpus. Das 46 referências com data elegível entre as
60 mais citadas, apenas 15 foram recuperadas pela busca (todas elas incluídas).
Ficaram de fora trabalhos que o próprio referencial teórico trata como
estruturantes --- \textcite{frey2017}, \textcite{acemogluRestrepo2020robots},
\textcite{webb2020}, \textcite{felten2021}, \textcite{eloundou2024},
\textcite{noyZhang2023} ---, seja por usarem vocabulário ausente da
\textit{string} (``\textit{computerisation}'', ``\textit{robots}''), seja por
circularem em periódicos de ciência geral ou como \textit{working papers}, fora
do alcance das bases. O corpus representa, assim, a literatura econômica
\emph{difusa} e indexada, e não a fronteira da disciplina. Isso não invalida a
pergunta da revisão, que é sobre o enquadramento predominante no campo, mas
delimita o que se pode concluir: os percentuais descrevem como a literatura
ampla absorve e reproduz os diagnósticos, e a fronteira entra na análise pelo
referencial teórico e pela rede de co-citação. Como os trabalhos de fronteira
ausentes sobre IA generativa apontam majoritariamente para a alta qualificação,
o viés de cobertura opera, se tanto, \emph{contra} a hipótese central.

\subsection{Síntese: três proposições}

A discussão permite condensar a resposta à pergunta da revisão em três
proposições. Primeira: o diagnóstico da literatura mudou com a IA generativa, e
mudou na direção prevista, mas a mudança acompanha o objeto tecnológico mais do
que o calendário --- a literatura sobre automação segue dizendo o que dizia, e a
literatura sobre IA generativa diz outra coisa. Segunda: essa diferença não
contradiz o arcabouço de tarefas, e sim o confirma, pois o risco segue o
conteúdo das tarefas que cada tecnologia executa; o que envelheceu não foi a
teoria, e sim a identificação entre ``rotineiro'' e ``automatizável''. Terceira:
a aparente divergência de vereditos sobre o emprego reflete sobretudo o que cada
desenho mede, e o campo caminha para saber mais sobre a firma e menos sobre o
agregado, justamente quando a pergunta relevante para a política pública é a
agregada.
